\documentclass[11pt,a4paper]{article}
\usepackage[utf8]{inputenc}
\usepackage[T1]{fontenc}
\usepackage[french]{babel}
\usepackage{amsmath,amssymb,amsthm}
\usepackage[margin=2.3cm]{geometry}
\usepackage{tikz}
\usetikzlibrary{decorations.pathreplacing,arrows.meta}
\usepackage[most]{tcolorbox}

\newtheorem*{prop}{Propriété}
\newcommand{\E}{\operatorname{E}}
\newcommand{\Q}{\mathbb{Q}}
\newcommand{\N}{\mathbb{N}}
\newcommand{\Z}{\mathbb{Z}}
\newcommand{\R}{\mathbb{R}}

\newtcolorbox{idee}{colback=blue!5,colframe=blue!50!black,title=Idée}
\newtcolorbox{attention}{colback=orange!8,colframe=orange!70!black,title=À retenir}

\title{Tout intervalle contient un rationnel\\[2pt]\large Démonstration détaillée avec la propriété d'Archimède}
\date{}

\begin{document}
\maketitle

\begin{prop}
Pour tous réels $a<b$, il existe $r\in\Q$ tel que $a<r<b$.
\end{prop}

\section{L'idée en une image}

Fixons un entier $q\ge 1$. Les fractions de dénominateur $q$, c'est-à-dire les $\dfrac{p}{q}$ avec $p\in\Z$, sont espacées de $\dfrac1q$ sur la droite réelle : ce sont les graduations d'une règle.

\begin{center}
\begin{tikzpicture}[scale=1.55, >=Stealth]
  % droite
  \draw[->] (-0.3,0) -- (4.2,0);
  % graduations p/q (pas 0.7)
  \foreach \k in {0,1,...,5}{
    \draw (0.7*\k,0.08) -- (0.7*\k,-0.08);
  }
  \node[below] at (0.0,-0.1) {\small $\frac{p-2}{q}$};
  \node[below] at (0.7,-0.1) {\small $\frac{p-1}{q}$};
  \node[below,red!70!black] at (1.4,-0.1) {\small $\boxed{\frac{p}{q}}$};
  \node[below] at (2.1,-0.1) {\small $\frac{p+1}{q}$};
  \node[below] at (2.8,-0.1) {\small $\frac{p+2}{q}$};
  \node[below] at (3.5,-0.1) {\small $\frac{p+3}{q}$};
  % a et b
  \draw[blue!60!black,very thick] (1.0,0.18) -- (1.0,-0.18);
  \draw[blue!60!black,very thick] (3.2,0.18) -- (3.2,-0.18);
  \node[above,blue!60!black] at (1.0,0.2) {$a$};
  \node[above,blue!60!black] at (3.2,0.2) {$b$};
  % intervalle
  \draw[blue!60!black,line width=2pt,opacity=.35] (1.0,0) -- (3.2,0);
  % point p/q
  \fill[red!70!black] (1.4,0) circle (2pt);
  % accolade b-a
  \draw[decorate,decoration={brace,amplitude=5pt}] (1.0,0.75) -- (3.2,0.75)
    node[midway,above=6pt] {\small longueur $b-a$};
  % accolade 1/q
  \draw[decorate,decoration={brace,amplitude=3pt,mirror}] (0.7,-0.85) -- (1.4,-0.85)
    node[midway,below=4pt] {\small pas $\frac1q$};
\end{tikzpicture}
\end{center}

\begin{idee}
Si le pas $\dfrac1q$ est \emph{strictement plus petit} que la longueur $b-a$ de l'intervalle, alors une graduation tombe forcément dans $]a,b[$ : on ne peut pas « sauter par-dessus » l'intervalle.
\end{idee}

\section{Étape 1 : choisir $q$ grâce à Archimède}

\begin{attention}
\textbf{Propriété d'Archimède.} Pour tout réel $M$, il existe un entier $q\in\N^*$ tel que $q>M$. Autrement dit, $\N$ n'est pas borné dans $\R$.
\end{attention}

Comme $a<b$, on a $b-a>0$, donc $M=\dfrac{1}{b-a}$ est bien un réel. Archimède donne alors $q\in\N^*$ tel que
\[
q>\frac{1}{b-a}.
\]
Multiplions par $b-a>0$ (le sens de l'inégalité est conservé) :
\[
q(b-a)>1
\quad\Longleftrightarrow\quad
\boxed{\frac1q<b-a}.
\]
C'est exactement la condition « pas $<$ longueur » du dessin.

\section{Étape 2 : choisir $p$ avec la partie entière}

On pose $p=\E(aq)+1$. Par définition de la partie entière, $\E(aq)\le aq<\E(aq)+1$, donc
\[
\boxed{p-1\le aq<p}.
\]
Ainsi $p$ est le plus petit entier \emph{strictement} supérieur à $aq$ : $\dfrac{p}{q}$ est la première graduation à droite de $a$.

\section{Étape 3 : vérifier $a<\dfrac pq<b$}

\paragraph{Minoration.} On divise $aq<p$ par $q>0$ :
\[
a<\frac pq.
\]

\paragraph{Majoration.} On divise $p-1\le aq$ par $q>0$ :
\[
\frac pq-\frac1q\le a
\quad\text{donc}\quad
\frac pq\le a+\frac1q.
\]
Or $\dfrac1q<b-a$ (étape 1), d'où
\[
\frac pq\le a+\frac1q<a+(b-a)=b.
\]

\section{Conclusion}

On a $a<\dfrac pq<b$ avec $p\in\Z$ et $q\in\N^*$, donc $r=\dfrac pq\in\Q$ convient. \hfill$\blacksquare$

\bigskip
\noindent\textbf{Résumé en trois lignes.}
\begin{enumerate}
  \item Archimède $\Rightarrow$ $q>\frac1{b-a}$, donc $\frac1q<b-a$ (le pas est plus petit que l'intervalle).
  \item $p=\E(aq)+1$ $\Rightarrow$ $p-1\le aq<p$ (première graduation après $a$).
  \item On en déduit $a<\frac pq\le a+\frac1q<b$.
\end{enumerate}

\end{document}
